\subsection{性交疼痛恐惧（Coitophobia）与初夜焦虑}

性交疼痛恐惧（coitophobia，又称 genophobia、性交恐惧症）指对性交产生强烈的、持续的恐惧与回避，程度超过实际疼痛本身，甚至一想到插入就出现心悸、出汗、肌肉僵硬。它与因疼痛而"不敢做"的合理反应不同，属于焦虑障碍的范畴，常与阴道痉挛互为因果。

\subsubsection{常见成因}

- \textbf{既往疼痛经历}：既往性交疼痛、妇科检查疼痛、外阴前庭痛，使大脑把"插入"与"疼痛"建立条件反射。
- \textbf{性创伤与暴力经历}：性侵、性骚扰、童年期不良经历是重要诱因。
- \textbf{错误信息与性羞耻}：相信"初夜必剧痛、必大量出血""性就是忍耐"等说法，把性交预设为痛苦与危险。
- \textbf{关系与心理因素}：对怀孕、失去控制、体像不满、亲密恐惧、伴侣压力等的深层焦虑。
- \textbf{阴道痉挛}：盆底肌不自主收缩使插入困难，疼痛与恐惧互相强化，形成恶性循环。

\subsubsection{与"初夜焦虑"的区别}

- \textbf{初夜焦虑}多是对"第一次"的普遍性紧张与期待，通常随经验增加而缓解。
- \textbf{性交疼痛恐惧}则是一种明显影响功能的持续恐惧，即使环境安全、身体放松也难以进行，需要干预。

\subsubsection{突破路径}

- \textbf{停止"硬闯"}：疼痛与恐惧状态下强行插入只会加深创伤，应先停下来。
- \textbf{逐步脱敏与自我探索}：借助镜子与手指了解身体，从非插入的亲密、外阴轻触开始，逐级过渡到手指、扩张器，每一步都在无痛、放松、可控的前提下进行。
- \textbf{盆底肌放松训练}：学习识别与放松盆底肌（反向凯格尔），配合腹式呼吸与渐进式肌肉放松。
- \textbf{认知行为治疗}：识别并修正"性=危险/疼痛"的自动化思维，建立安全预期。
- \textbf{伴侣参与}：伴侣的理解、耐心与"不进则止"的配合是治疗成功的关键；把目标从"完成插入"改为"共同感到舒适"。
- \textbf{专业治疗}：性治疗师、盆底物理治疗师与心理科医生的多学科协作；合并阴道痉挛者可参照上文扩张器方案。

\noindent\textit{【编者按】}若恐惧已严重影响性生活、情绪或自我认同，请及时寻求专业帮助——这不是"矫情"或"不够爱对方"，而是可以治疗的医学问题。多数人可在数周到数月内改善，急于求成反而延长病程。
